\documentclass[11pt,a4paper]{article}

\usepackage[utf8]{inputenc}
\usepackage[T1]{fontenc}
\usepackage[left=2cm,right=2cm,top=2cm,bottom=2cm]{geometry}
\usepackage{hyperref}
\usepackage{minted}
\usepackage[svgnames]{xcolor}
\usepackage{amsmath,amsfonts,amssymb}

\usepackage{minifeyn}

\newcommand{\testprop}[1]{%
    \begin{inlinefeyn}
        \draw[#1] (0,0) -- (0.5,0);
    \end{inlinefeyn}
}

\author{Killian Bouzoud}
\title{The \texttt{minifeyn} package}
\date{\texttt{v1.1.0}}

\begin{document}

\maketitle

\begin{abstract}
    \texttt{minifeyn} is a package to draw Feynman diagrams in equations.
    It uses \texttt{TikZ} and allows for customizing the diagrams and working with delimiters.
\end{abstract}

\tableofcontents

\newpage

\section{``Install'' the package}

Simply put the \texttt{minifeyn.sty} file in the same directory as your document and load it using:
\begin{minted}[bgcolor=Beige]{latex}
    \usepackage{minifeyn}
\end{minted}
That's pretty much it.

\subsection{Requirements}

\texttt{minifeyn} uses the following \LaTeX~packages:
\begin{itemize}
    \item \texttt{boolexpr}
    \item \texttt{TikZ} and the following libraries:
    \begin{itemize}
        \item \texttt{decorations.pathmorphing}
        \item \texttt{decorations.markings}
        \item \texttt{arrows.meta}
    \end{itemize}
\end{itemize}

\section{Using the package}

\subsection{Available particle styles}
\label{sec:particlestyles}

The following particles can be used in diagrams:
\begin{alignat}{3}
    \testprop{mnf-sca}\qquad     && &\texttt{mnf-sca}\qquad     && \text{Scalar}\nonumber                                   \\
    \testprop{mnf-chasca}\qquad  && &\texttt{mnf-chasca}\qquad  && \text{Charged scalar}\nonumber                           \\
    \testprop{mnf-fer}\qquad     && &\texttt{mnf-fer}\qquad     && \text{Fermion}\nonumber                                  \\
    \testprop{mnf-dfer}\qquad    && &\texttt{mnf-dfer}\qquad    && \text{Double fermion (spin-$3/2$ field, \dots)}\nonumber \\
    \testprop{mnf-bos}\qquad     && &\texttt{mnf-bos}\qquad     && \text{Boson}\nonumber                                    \\
    \testprop{mnf-dbos}\qquad    && &\texttt{mnf-dbos}\qquad    && \text{Double boson (spin-$2$ field, \dots)}\nonumber     \\
    \testprop{mnf-gaugino}\qquad && &\texttt{mnf-gaugino}\qquad && \text{Gaugino}
\end{alignat}

\subsection{Available diagrams}
\label{sec:diagrams}

\subsubsection{One-loop diagrams}

\begin{minted}[bgcolor=Beige]{latex}
    \mnfoneloop[diag]{s1}{s2}
\end{minted}
There are two topologies:
\begin{alignat}{2}
    &\begin{inlinefeyn}
        \mnfoneloop[B]{mnf-fer}{mnf-fer}
    \end{inlinefeyn}\qquad
    &&
    \begin{inlinefeyn}
        \mnfoneloop[T]{mnf-fer}{mnf-fer}
    \end{inlinefeyn}\nonumber \\
    &\texttt{B}\qquad && \texttt{T}
\end{alignat}
The topology can be selected by the (optional) argument \texttt{diag}, which can either be \texttt{B} (bubble, \textbf{default}) or \texttt{T} (tadpole).
The \texttt{s1} and \texttt{s2} arguments are the particle styles of the external line and loop, respectively.
They can be any of the styles in Section~\ref{sec:particlestyles}.

\subsubsection{$1\leftrightarrow 2$ diagram}

\begin{minted}[bgcolor=Beige]{latex}
    \mnfonetotwo{s1}{s2}{s3}
\end{minted}
There is only one topology:
\begin{equation}
    \begin{inlinefeyn}
        \mnfonetotwo{mnf-fer}{mnf-fer}{mnf-fer}
    \end{inlinefeyn}
\end{equation}
The \texttt{s1}, \texttt{s2} and \texttt{s3} arguments are the styles of each particle, ordered from left to right and top to bottom.

\subsubsection{$2\leftrightarrow 2$ diagrams}

\begin{minted}[bgcolor=Beige]{latex}
    \mnftwototwo[diag]{s1}{s2}{s3}{s4}{s5}
\end{minted}
There are three topologies:
\begin{alignat}{3}
    &\begin{inlinefeyn}
        \mnftwototwo[S]{mnf-fer}{mnf-fer}{mnf-fer}{mnf-fer}{mnf-fer}
    \end{inlinefeyn}\qquad
    &&
    \begin{inlinefeyn}
        \mnftwototwo[T]{mnf-fer}{mnf-fer}{mnf-fer}{mnf-fer}{mnf-fer}
    \end{inlinefeyn}\qquad
    &&
    \begin{inlinefeyn}
        \mnftwototwo[X]{mnf-fer}{mnf-fer}{}{mnf-fer}{mnf-fer}
    \end{inlinefeyn}\nonumber \\
    &\texttt{S}\qquad && \texttt{T}\qquad && \texttt{X}
\end{alignat}
The topology can be selected by the (optional) argument \texttt{diag},
which can either be \texttt{S} ($s$ channel), \texttt{T} ($t$ channel, \textbf{default}) or \texttt{X} (four-point interaction).
The \texttt{s1}, \texttt{s2}, \texttt{s4} and \texttt{s5} arguments are the styles of each particle, ordered from top to bottom and left to right.
The \texttt{s3} argument is the style of the intermediate particle, which does nothing for \texttt{X}.

For any of the functions that have a \texttt{diag} optional argument, if the value is not recognized the following will be outputted instead:
\begin{equation}
    \begin{inlinefeyn}
        \mnfoneloop[E]{}{}
    \end{inlinefeyn}
\end{equation}

\subsection{Writing equations}

The functions from Section~\ref{sec:diagrams} should be used wrapped in a \texttt{inlinefeyn} environment.
For example, the following code:
\inputminted[bgcolor=Beige]{latex}{include/equation1.tex}
produces the following output
\begin{equation}
    \begin{inlinefeyn}
        \mnftwototwo{mnf-chasca}{mnf-chasca}{mnf-bos}{mnf-chasca}{mnf-chasca}
    \end{inlinefeyn}
\end{equation}
The \texttt{minifeyn} functions are meant be used in equations.
For example, we can write the diagrams for Bhabha scattering as
\inputminted[bgcolor=Beige]{latex}{include/equation2.tex}
which gives us
\begin{equation}
    i\mathcal{M}^{ }_{\rm Bhabha}=
    \begin{inlinefeyn}
        \mnftwototwo{mnf-fer}{mnf-fer}{mnf-bos}{mnf-fer}{mnf-fer}
    \end{inlinefeyn}
    +
    \begin{inlinefeyn}
        \mnftwototwo[S]{mnf-fer}{mnf-fer}{mnf-bos}{mnf-fer}{mnf-fer}
    \end{inlinefeyn}
\end{equation}
Delimiters scale properly with the Section~\ref{sec:diagrams} functions.
For example, we can write the \textit{square} of the previous expression
\begin{equation}
    |\mathcal{M}^{ }_{\rm Bhabha}|^2=
    \left|\begin{inlinefeyn}
        \mnftwototwo{mnf-fer}{mnf-fer}{mnf-bos}{mnf-fer}{mnf-fer}
    \end{inlinefeyn}\right|^2
    +\left|\begin{inlinefeyn}
        \mnftwototwo[S]{mnf-fer}{mnf-fer}{mnf-bos}{mnf-fer}{mnf-fer}
    \end{inlinefeyn}\right|^2
    +\begin{inlinefeyn}
        \mnftwototwo{mnf-fer}{mnf-fer}{mnf-bos}{mnf-fer}{mnf-fer}
    \end{inlinefeyn}\times\left(\begin{inlinefeyn}
        \mnftwototwo[S]{mnf-fer}{mnf-fer}{mnf-bos}{mnf-fer}{mnf-fer}
    \end{inlinefeyn}\right)^{\star}
    +\begin{inlinefeyn}
        \mnftwototwo[S]{mnf-fer}{mnf-fer}{mnf-bos}{mnf-fer}{mnf-fer}
    \end{inlinefeyn}\times\left(\begin{inlinefeyn}
        \mnftwototwo{mnf-fer}{mnf-fer}{mnf-bos}{mnf-fer}{mnf-fer}
    \end{inlinefeyn}\right)^{\star}
\end{equation}
as
\inputminted[bgcolor=Beige]{latex}{include/equation3.tex}
Note that the \texttt{inlinefeyn} environment is actually just a \texttt{tikzpicture}
(which also adjusts the vertical alignment and sets the width of the lines in the figure),
meaning it is possible to complement the diagrams with other \texttt{TikZ} instructions.
For example, we could start from the ``bubble'' diagram and add a line to get the following two-loop diagram for the $\phi^4$ model
\begin{equation}
    \begin{inlinefeyn}
        \mnfoneloop{mnf-sca}{mnf-sca}
        \draw[mnf-sca] (0.5,0) -- (1.5,0);
    \end{inlinefeyn}
\end{equation}
by writing
\inputminted[bgcolor=Beige]{latex}{include/equation4.tex}
The coordinates are chosen such that:
\begin{itemize}
    \item All diagrams are $1$ unit high,
    \item The $y$ range is $[0,1]$ for \texttt{mnftwototwo} and \texttt{mnfoneloop[T]} and $[-0.5,0.5]$ for all other diagrams,
    \item The $x$ range is $[0,2]$ for loop diagrams and $[0,1]$ everywhere else,
    \item The $(0,0)$ point is the leftmost point, for $2\leftrightarrow 2$ diagrams it is at the bottom of the diagram.
\end{itemize}
You can see the source code for more information.

\section{Examples}

Some more examples are listed below.

\subsection{Axion production rate}

\begin{align}
    \Gamma^{ }_\varphi\Big|^{ }_{2\leftrightarrow 2}\equiv\frac{1}{4k n^{ }_{\rm B}}
    \int {\rm d}\Omega^{ }_{2\leftrightarrow 2}&\left\{\left|\begin{inlinefeyn}
        \mnftwototwo{mnf-fer}{mnf-bos}{mnf-bos}{mnf-fer}{mnf-sca}
    \end{inlinefeyn}\right|^2\times f^{ }_q f^{ }_g\left[1-f^{ }_q\right]
    \right.\nonumber \\
    &+\left.\left|\begin{inlinefeyn}
        \mnftwototwo{mnf-bos}{mnf-bos}{mnf-bos}{mnf-bos}{mnf-sca}
    \end{inlinefeyn}\right|^2\times f^{ }_g f^{ }_g\left[1+f^{ }_g\right]
    +\cdots\right\}
\end{align}
\inputminted[bgcolor=Beige]{latex}{include/equation5.tex}

\subsection{Cutting rules}

\begin{equation}
    \begin{inlinefeyn}
        \mnfoneloop{mnf-bos}{mnf-fer}
        \draw[mnf-bos] (0.646447,-0.353553) -- (1.35355,0.353553);
        \draw[mnf-fer,red] (0.8,-0.6) -- (1.0,-0.6) -- (1.0,0.6) -- (1.2,0.6);
    \end{inlinefeyn}\to\begin{inlinefeyn}
        \mnfonetotwo{mnf-bos}{mnf-fer}{mnf-fer}
        \draw[mnf-bos] (0.75,-0.25) -- (1.0,0.0);
    \end{inlinefeyn}+\begin{inlinefeyn}
        \mnftwototwo[S]{mnf-fer}{mnf-bos}{mnf-fer}{mnf-bos}{mnf-fer}
    \end{inlinefeyn}+\cdots
\end{equation}
\inputminted[bgcolor=Beige]{latex}{include/equation6.tex}

\section{Known issues}

\subsection{Interaction with unknown package}

When writing my PhD thesis and using a preliminary version of this package,
I ended up running into a bug where I couldn't place two diagrams on the same line unless I wrapped one of the \texttt{inlinefeyn} environments
in a \mintinline{latex}{\raisebox{0ex}{...}}.
I have no clue what caused this and I did not run into this bug again while developping the package in this repo.
My thesis uses many more \LaTeX~packages than this manual though, so the issue is probably caused by that?
Again, I do not know the cause or if you might even encounter that bug but I think it's useful to have that information.

\subsection{Interaction with \texttt{annotate-equations}}

When using the \texttt{annotate-equations} package in conjunction with \texttt{minifeyn},
the \mintinline{latex}{\eqnmarkbox} command can be used to put a colored box around a \mintinline{latex}{\inlinefeyn} environment
(like with delimiters, the size is adjusted to the diagram).
\textbf{However} when trying to place an annotation arrow with \mintinline{latex}{\annotate}, the arrow will not be drawn.
I do not know what causes this either.

\end{document}
